% Milestone 2b: experiment tracking. Included from sections/2-2-reproducibility.tex,
% so it appears as subsection 2.2's last sub-subsection.

\subsubsection{Experiment tracking}
\label{sec:m2-experiment-tracking}
\guidance{
  \textbf{Template:} e.g.\ MLflow.\\
  \textbf{Rubric (5 pts):} How does our experiment tracking help us and others compare, understand, and extend our results?\\
  The demo uses DagsHub as shared MLflow tracking server with credentials in \texttt{.env}.
  Show what we log (parameters, metrics, artifacts, emissions) and a comparison of runs.
  Our own starting points: one experiment \texttt{recommenditos-price} (\texttt{train.mlflow\_experiment} in \texttt{params.yaml}), one run per variant of the experiment ladder (\texttt{b0}, \texttt{b1}, \texttt{lgbm-basic}, \texttt{lgbm-extended}), and the DagsHub repository of EDN-19 as tracking server.
}

We track experiments with MLflow on the team's DagsHub repository, which is also the DVC remote (EDN-19).
The credentials are in a \texttt{.env} file that Git ignores.

\paragraph{What we track.}
The experiment \texttt{recommenditos-price} holds one run per variant of our experiment ladder (\cref{tab:runs}), a sequence of models from a simple baseline to the full model.
The ladder is defined in \texttt{params.yaml} and has four steps: \texttt{b0}, which predicts the median training price of cars with the same make, model, and two-year age band; \texttt{b1}, a Ridge regression; \texttt{lgbm-basic}, LightGBM on the 16 basic features, trained on its rows plus blanked copies of them (EDN-82); and \texttt{lgbm-extended}, LightGBM without such copies on 162 features, which add the equipment lists, the history flags, and further technical columns.
Each run logs:

\begin{itemize}
  \item the parameters of its stage and the name of its feature set;
  \item the metrics: MdAPE, the shares within $\pm$10\,\% and $\pm$20\,\%, the mean absolute error in EUR, and the success criteria;
  \item the fit time and CodeCarbon's energy and emissions figures (\cref{sec:m3-energy});
  \item the model artefact;
  \item the Git commit and the hashes of the stage's inputs, so that a run can be matched to a committed state of the pipeline (NFR-14).
\end{itemize}

\paragraph{How others can compare and extend the results.}
Each variant is its own DVC stage, so \texttt{dvc.lock} pins the code, data, and parameters behind a run, and \texttt{dvc repro train@lgbm-basic} on the same commit reproduces it.
The comparison is also committed: the \texttt{evaluate} stage writes one record per variant to \texttt{reports/metrics/} and a summary to \texttt{metrics.json}.
A pull request that changes a result shows the change in its diff, and a reader needs no account on the tracking server.
A new variant is one more entry in \texttt{params.yaml}.

\begin{table}[ht]
  \centering\small
  \begin{tabular}{@{}lrrrrrr@{}}
    \toprule
    & & \multicolumn{3}{c}{Every field given} & Required fields only & \\
    \cmidrule(lr){3-5}\cmidrule(lr){6-6}
    Run & Features & MdAPE & Within $\pm$20\,\% & MAE (EUR) & MdAPE & Fit time \\
    \midrule
    \texttt{b0} & 16 & 12.16\,\% & 70.5\,\% & 9,388 & 12.16\,\% & 0.1\,s \\
    \texttt{b1} & 16 & 9.52\,\% & 80.3\,\% & 7,499 & 10.72\,\% & 2.0\,s \\
    \texttt{lgbm-basic} & 16 & 6.73\,\% & 90.0\,\% & 4,731 & 6.86\,\% & 19.0\,s \\
    \texttt{lgbm-extended} & 162 & 6.26\,\% & 90.4\,\% & 4,651 & 10.30\,\% & 30.5\,s \\
    \bottomrule
  \end{tabular}
  \caption{The experiment ladder on the 19,665 test rows, from \texttt{reports/metrics/} and \texttt{reports/figures/energy-vs-error.csv}. The MLflow run ID of each row is in \texttt{reports/metrics/<run>.json}. Runs of 2026-10-13 on a laptop CPU.}
  \label{tab:runs}
\end{table}

\paragraph{What the comparison showed.}
\Cref{tab:runs} compares the four runs, which the \href{https://dagshub.com/recommenditos/MLOPS_Recommenditos/experiments}{experiment on DagsHub} lists as well.
The extended feature set improves the MdAPE by 0.48 percentage points when the user gives every field.
With only the fields FR-01 requires, \texttt{lgbm-extended} is at 10.30\,\%, which is 1.65 times its own full-input value and above SC-06's bound of 1.5, and it is worse than \texttt{lgbm-basic} at 6.86\,\%.
Such requests were almost absent from the training rows, and \texttt{lgbm-basic} was at 8.26\,\% on them until we trained it on blanked copies (EDN-82).
\texttt{lgbm-extended} is trained without such copies.
With them, its error on a request of the required fields fell to 7.14\,\% on the validation split, but it then lost to \texttt{lgbm-basic} on a request that gives exactly the six optional fields of the basic set.
\texttt{lgbm-extended} also takes 1.6 times as long to fit.
We therefore put \texttt{lgbm-basic} forward as the candidate (EDN-62).

The hyperparameter search is tracked separately.
Its 34 fits are runs in a second experiment, \texttt{recommenditos-price-tuning}, so they do not mix with the ladder.
The search compared \texttt{learning\_rate}, \texttt{num\_leaves}, and \texttt{min\_child\_samples} on the validation split for both LightGBM variants and found no value that is measurably better than the committed ones.
The results in \cref{tab:runs} use these values; the search was not repeated after \texttt{lgbm-basic} got its copies.
